\documentclass[10pt,a4paper]{article}

% ---------------------------------------------------------------------------
% Packages
% ---------------------------------------------------------------------------
\usepackage[margin=2.1cm]{geometry}
\usepackage[T1]{fontenc}
\usepackage{lmodern}
\usepackage{microtype}
\usepackage{amsmath,amssymb}
\usepackage{booktabs}
\usepackage{tabularx}
\usepackage{multirow}
\usepackage{graphicx}
\usepackage{xcolor}
\usepackage{siunitx}
\usepackage{enumitem}
\usepackage{caption}
\usepackage[numbers,sort&compress]{natbib}
\usepackage[colorlinks=true,linkcolor=blue!50!black,citecolor=blue!50!black,urlcolor=blue!50!black]{hyperref}

\sisetup{group-separator={,},group-minimum-digits=4}
\setlist{nosep,leftmargin=*}
\captionsetup{font=small,labelfont=bf}
\setlength{\parskip}{0.3em}

% Placeholder for numbers that the full suite has not produced yet. Remove the
% definition's colour once every \TBD{} is filled in; `grep -n TBD report.tex`
% lists what is outstanding.
\newcommand{\TBD}[1]{\textcolor{red}{\textbf{[#1]}}}

% Figure that renders a framed placeholder until the file exists in figures/.
\newcommand{\resultfigure}[3]{%
  \IfFileExists{figures/#1}%
    {\includegraphics[width=#2]{figures/#1}}%
    {\fbox{\parbox[c][0.17\textheight][c]{#2}{\centering\textcolor{red}{\textbf{[Figure placeholder: \detokenize{#1}]}}\\[0.4em]\small #3}}}%
}

\newcommand{\BO}{\textsc{B0}}
\newcommand{\GQA}[1]{GQA-#1}

\title{\vspace{-1.2em}Grouped-Query Attention and Partial RoPE Under a Fixed FLOP Budget:\\
A Controlled Study of a Small Language Model Trained from Scratch}
\author{Jerry Yang Shao Ern (A0281549U)}
\date{}

\begin{document}
\maketitle
\vspace{-2.2em}

% ---------------------------------------------------------------------------
% Declarations (required by the brief; stated before any results)
% ---------------------------------------------------------------------------
\begin{center}
\fbox{\parbox{0.96\linewidth}{\small
\textbf{Repository:} \url{https://github.com/mcxraider/modded-nanochat}
(experiment namespace \texttt{d8-v5}; code commit \TBD{commit hash of the d8-v5 launch}).\\
\textbf{Group members:} Jerry Yang Shao Ern, student number A0281549U.\\[0.3em]
\textbf{Compute budget (declared before any training run).} Hardware: one NVIDIA H100 SXM 80\,GB on Modal, BF16 autocast,
FlashAttention-3. Per-run allowance: at most $C = 1.385\times10^{17}$ \emph{estimated} training FLOPs, the cost of the
baseline \BO{} trained on upstream nanochat's compute-optimal horizon at depth 8 (1{,}680 updates $\times$ 262{,}144
tokens $\times$ 314{,}575{,}872 FLOPs/token). FLOPs are counted with nanochat's estimator (6 FLOPs per matmul parameter per
token plus windowed attention), so each architecture receives the largest whole number of updates whose estimate fits in
$C$. A 2-hour wall-clock cap applies to every run as a safety stop only; no comparable run reached it (each needed about
\TBD{6--9} minutes). Pilots, numerical preflight checks and tokenizer training are excluded from the budget and reported
separately.\\[0.3em]
\textbf{Initialisation.} Every model is trained from random initialisation. No pretrained weights, embeddings or other
pretrained neural components are used. We reuse the nanochat architecture and code, its BPE tokenizer recipe (trained
by us on the training shards, vocabulary 32{,}768) and the ClimbMix corpus.\\[0.3em]
\textbf{Evaluation objective.} Bits per byte (BPB) of next-token prediction on held-out ClimbMix web text: $2^{24}$ tokens
from shard 06542 (validation, used for curves and any decision) and, separately, $2^{24}$ tokens from shard 06541 (test,
evaluated once per run after training; the primary metric). BPB normalises the summed cross-entropy by the UTF-8 byte
count of the targets, so it is independent of tokenisation and measures the quantity the budget is spent on, namely
language modelling of the training distribution. It does not measure downstream reasoning, instruction following or
long-context use, and inference speed and memory are reported as separate metrics.}}
\end{center}

% ---------------------------------------------------------------------------
\section{Introduction and research question}
% ---------------------------------------------------------------------------
Two attention changes appear in almost every recent open model: \emph{grouped-query attention}
(GQA)~\cite{ainslie2023gqa,shazeer2019mqa}, which lets several query heads share one key/value (K/V) head, and
\emph{partial rotary embeddings}~\cite{su2021roformer,black2022gptneox}, which rotate only a fraction of each head's
dimensions. Both are usually justified by inference cost (a smaller K/V cache) or by evidence from large models, not by
what they do to a small model trained from scratch on a tight compute budget. The brief asks how to obtain the best
model within a fixed per-run allowance, which makes this the relevant question. GQA removes training FLOPs, so at fixed
compute a GQA model trains on \emph{more} tokens than its multi-head (MHA) counterpart. Whether it ends up better or
worse therefore depends on two competing effects: a per-token quality cost from sharing K/V heads, and a gain from the
extra data. Our research question is:

\begin{quote}\itshape
At depth 8 and a fixed per-run budget of $C$ estimated training FLOPs, what do GQA and partial RoPE cost or buy in
held-out BPB, training throughput and inference memory/speed, and how much of GQA's fixed-compute result comes from
the extra tokens its FLOP saving buys?
\end{quote}

We wrote the hypotheses, resolution thresholds and analysis rules in \texttt{experiments/experiments.md} before running
any training job. Every outcome below is reported against its prediction, including the null and negative ones.
Our contributions are: (i) a fixed-compute comparison of MHA, GQA-2 and MQA with a seed-noise estimate on
\emph{both} the baseline and the main variant; (ii) a dose--response over the rotary fraction; (iii) two controls for
confounds that the naive comparison hides: the \emph{value embeddings} in this architecture shrink with the K/V head
count (a $2\times2$ design), and GQA's per-token penalty is mixed with its extra data (a matched-token ablation);
and (iv) an inference benchmark that tests where the K/V saving actually appears on an H100.

% ---------------------------------------------------------------------------
\section{Background}
% ---------------------------------------------------------------------------
\paragraph{Multi-query and grouped-query attention.}
In MHA~\cite{vaswani2017attention} every query head $h$ has its own key and value projections. During autoregressive
decoding, each step must read the whole cache of past keys and values, so at large batch or long context decoding is
limited by memory bandwidth rather than arithmetic~\cite{shazeer2019mqa,pope2023efficiently}.
MQA~\cite{shazeer2019mqa} shares a single K/V head across all query heads and cuts the cache by a factor of $n_\text{head}$.
GQA~\cite{ainslie2023gqa} interpolates between the two with $n_\text{kv}$ groups, and reports quality close to MHA when
large checkpoints are \emph{uptrained} into GQA form. That evidence comes from converted T5-XXL models, not from small
models trained from scratch, which is why we test the claim directly. In training, GQA saves only the K/V projection
matmuls: the attention score and value products still run over all $n_\text{head}$ query heads. At our size this is
$8\%$ of B0's estimated FLOPs for full K/V, so GQA-2 saves $4.2\%$ and MQA $6.4\%$ (Table~\ref{tab:arch}).

\paragraph{Rotary embeddings and partial RoPE.}
RoPE~\cite{su2021roformer} rotates pairs of query/key dimensions by angles proportional to position, so the dot product
depends on relative offset. Each of the $d/2$ pairs has frequency $\theta_i = 10000^{-2i/d}$. GPT-NeoX~\cite{black2022gptneox}
rotates only 25\% of each head and leaves the rest position-free. Barbero et al.~\cite{barbero2025round} found that a
trained Gemma model uses the high-frequency pairs for positional patterns and the lowest frequencies as almost
position-independent ``semantic'' channels, and that removing the lowest frequencies can help. Partial RoPE achieves
something similar by construction. A detail matters for interpreting our dose--response: in our implementation the
frequencies are recomputed for the rotated prefix (base $10^4$ over $d_\text{rot}$ dimensions), so R25 keeps the same
highest and lowest frequency as full RoPE but samples the spectrum $4\times$ more coarsely (16 pairs instead of 64), and
adds 96 dimensions with no positional signal at all.

\paragraph{Value embeddings.}
nanochat's attention adds a learned, token-indexed \emph{value embedding} to the values in alternating layers, scaled
by an input-dependent gate. This is a form of ResFormer value-residual learning~\cite{zhou2024value}, which reports
the same loss with 16\% fewer parameters. The embedding width equals the value width, $n_\text{kv}\cdot d_\text{head}$.
At depth 8 these tables hold 67.1M of B0's 125.8M parameters, so halving the K/V heads removes 33.5M embedding
parameters, about 16 times the 2.1M removed from the K/V projections. Nanochat's FLOP estimator treats embedding lookups as free,
so this change is invisible to the budget. A naive ``GQA vs MHA'' comparison in this codebase is therefore also a
``half vs full value embeddings'' comparison, and Experiment E5 separates the two.

\paragraph{Compute-optimal training.}
Chinchilla~\cite{hoffmann2022chinchilla} models loss as $L(N,D)=E+A/N^{\alpha}+B/D^{\beta}$ and finds roughly 20 tokens
per parameter compute-optimal. Kaplan et al.~\cite{kaplan2020scaling} first pointed out that embedding parameters
behave differently in such fits. nanochat~\cite{karpathy2025nanochat} fits its own ratio of 10.5 tokens per
\emph{scaling} parameter (transformer matrices plus unembedding) and we adopt its horizon for B0. We use the Chinchilla
data term only to make an order-of-magnitude prediction of what GQA's extra 4--6\% of tokens should be worth (\S\ref{sec:design}).

% ---------------------------------------------------------------------------
\section{Methods and experimental setup}
% ---------------------------------------------------------------------------
\subsection{Model, data and training procedure}
\paragraph{Baseline \BO.}
B0 is the nanochat ClimbMix leaderboard model at upstream commit \texttt{324e69c} (the configuration immediately before
the ``autoresearch'' rounds), rescaled to depth 8: 8 layers, width 512, 4 heads of dimension 128, context 2{,}048,
vocabulary 32{,}768. It keeps that commit's design: pre-norm with parameter-free RMSNorm, QK-norm~\cite{henry2020qknorm},
squared-ReLU MLP~\cite{so2021primer} of width $4\times$, untied embeddings, logit soft-capping, learnable residual and
input-skip scalars, gated value embeddings in alternating layers, and sliding-window attention with the
pattern SSSL (short windows of 1{,}024 tokens, every fourth and the final layer full)~\cite{beltagy2020longformer}.
We made one deliberate change: we train in BF16 rather than the leaderboard's FP8, so that numerical precision cannot
interact with the head shapes.

\paragraph{Variants.}
We extended the model with three switches, each verified by CPU unit tests and a GPU preflight
(\S\ref{sec:validity}): the number of K/V heads $n_\text{kv}\in\{4,2,1\}$; the rotary fraction
$f\in\{1, 0.75, 0.25\}$, where the first $f\cdot d_\text{head}$ dimensions of $q$ and $k$ are rotated and the rest pass
through unchanged; and the MLP ratio and value-embedding on/off flags used by the controls.
Table~\ref{tab:arch} lists the eight resulting architectures.

\paragraph{Data and tokenizer.}
We use NVIDIA ClimbMix~\cite{diao2025climb}, a 400B-token web mixture produced by clustering-based mixture search, in
the shard format used by nanochat. Upstream switched to ClimbMix because it gave clear gains over FineWeb-Edu and other
corpora on its leaderboard. The training pool is shards 0--47. Shard 06542 (the last, which nanochat treats as validation)
is our validation split, and shard 06541 is a held-out test split that is never used for any decision. We trained the
nanochat BPE tokenizer (vocabulary $2^{15}$) ourselves. Its trainer reads files in index order and stops after
$2\times10^9$ characters, so it sees only the first training shards and never the test shard \TBD{confirm from the
\texttt{tok\_train} log}. Documents are packed into 2{,}048-token rows with
nanochat's BOS-aligned best-fit packer. The longest run consumes \TBD{x.xx}\% of the pool, so no run repeats data
(\texttt{data\_epoch}$\,=1$ is checked).

\paragraph{Optimisation.}
Following the pinned upstream training script, matrices are updated with Muon~\cite{jordan2024muon} (Polar-Express
orthogonalisation~\cite{amsel2025polar} and NorMuon-style factored second moments~\cite{li2025normuon}) and the
embeddings, unembedding and scalars with AdamW. Upstream's own depth-8 rules set the global batch to 262{,}144
tokens, all learning rates scaled by $\sqrt{262144/524288}=0.707$, and the Muon weight decay to $0.371$, computed once from
B0 and \emph{frozen for every variant}. The schedule has no warm-up, a constant learning rate for 50\% of updates and
then linear decay to zero. Weight decay decays linearly with progress, and Muon momentum warms up from 0.85 to 0.95
over 300 updates. We tuned no hyperparameters for any variant (\S\ref{sec:limits}).

\subsection{The FLOP budget and why we chose it}\label{sec:budget}
We chose a FLOP budget rather than wall-clock time because (a) it is hardware-independent and reproducible, while
throughput on a rented GPU varies with the node; (b) it is the quantity scaling laws reason about; and (c) it rewards an
architecture exactly for the arithmetic it removes, which is the claimed advantage of GQA in training. The cost is that
the estimator ignores work that is not a matmul, so it understates the true cost of value embeddings and treats
partial RoPE as free. We report measured throughput next to every comparison so that readers can see how the ranking
would change under a time budget (\S\ref{sec:speed}). Training stops at the token target and refuses to start if the
target exceeds $C$. All fixed-compute runs use at least $99.98\%$ of $C$.

\begin{table}[t]
\centering\small
\caption{The eight architectures (depth 8, width 512, 4 query heads of dimension 128). FLOPs/token is nanochat's
estimate. ``Updates'' is the largest whole number of 262{,}144-token updates within $C$. Scaling parameters are the
transformer matrices plus the unembedding, and total parameters add the token and value embeddings. KV cache is in bytes per token
(BF16, all layers).}
\label{tab:arch}
\setlength{\tabcolsep}{4.5pt}
\begin{tabular}{@{}llrrrrr@{}}
\toprule
Name & Change from B0 & FLOPs/token & Updates & Tokens & Scaling / total params & KV B/token \\
\midrule
B0 & --- (MHA, $n_\text{kv}=4$) & 314.58M & 1{,}680 & 0.440B & 41.94M / 125.83M & 16{,}384 \\
GQA-2 & $n_\text{kv}=2$ & 301.99M & 1{,}750 & 0.459B & 39.85M / 90.18M & 8{,}192 \\
GQA-2-R75 & GQA-2, 75\% RoPE & 301.99M & 1{,}750 & 0.459B & 39.85M / 90.18M & 8{,}192 \\
GQA-2-R25 & GQA-2, 25\% RoPE & 301.99M & 1{,}750 & 0.459B & 39.85M / 90.18M & 8{,}192 \\
GQA-2-PM & GQA-2, MLP $4.5\times$ & 314.57M & 1{,}680 & 0.440B & 41.94M / 92.27M & 8{,}192 \\
MQA (GQA-4) & $n_\text{kv}=1$ & 295.70M & 1{,}787 & 0.468B & 38.80M / 72.35M & 4{,}096 \\
B0-noVE & B0, no value emb. & 314.57M & 1{,}680 & 0.440B & 41.94M / 58.72M & 16{,}384 \\
GQA-2-noVE & GQA-2, no value emb. & 301.99M & 1{,}750 & 0.459B & 39.85M / 56.62M & 8{,}192 \\
\bottomrule
\end{tabular}
\end{table}

\subsection{Experimental design}\label{sec:design}
We ran 14 training runs. Apart from the stated change, all share data, tokenizer, context, batch, optimiser, schedule and evaluation text, and all variants use seed 1 unless noted. The seed sets both the initialisation and the training-shard rotation,
so seed noise includes data-order noise.
\begin{itemize}
\item \textbf{E1, noise floor} (\texttt{d8-b0-s1..s3}): B0 with seeds 1--3. The sample s.d.\ $s$ of test BPB defines the resolution threshold.
\item \textbf{E2, main comparison} (\texttt{d8-gqa2-s1..s3}, \texttt{-r75}): GQA-2 (seeds 1--3) against the B0 mean, and GQA-2-R75 against both. The three GQA-2 seeds also check that
  the noise is similar across architectures.
\item \textbf{E3, parameter-matched control} (\texttt{d8-gqa2-pm}): GQA-2-PM spends GQA's saving on a wider MLP (2{,}304 instead of 2{,}048)
  instead of extra tokens. It matches B0's scaling parameters to 256 parameters and its FLOPs to $0.0005\%$.
\item \textbf{E4, dose--response} (\texttt{d8-gqa4}, \texttt{d8-gqa2-r25}): MQA (K/V axis, against B0 and GQA-2) and GQA-2-R25 (rotary axis, against GQA-2 and R75).
\item \textbf{E5, value-embedding $2\times2$} (\texttt{d8-b0-nove}, \texttt{d8-gqa2-nove}): B0-noVE and GQA-2-noVE. With B0 and GQA-2 these cross
  $\{4,2\}$ K/V heads with value embeddings $\{\text{on},\text{off}\}$.
\item \textbf{E6, matched-token ablation} (\texttt{d8-gqa2-tok}, \texttt{d8-gqa4-tok}): GQA-2 and MQA trained on B0's 1{,}680 updates (96\% and 94\% of $C$). The
  difference to their fixed-compute twin (same seed, so the same initial weights and batches up to update 1{,}680) is the
  value of the extra tokens, and the difference to B0 is the per-token cost of K/V sharing.
\end{itemize}

\paragraph{Pre-registered predictions.}
We wrote these predictions down before any run. \textbf{H2a:} GQA-2 is not resolved-worse than B0 at fixed compute.
\textbf{H2b:} GQA-2's training-throughput gain is below 5\%. \textbf{H2c:} R75 is not resolved from full RoPE in BPB
and differs by under 2\% in throughput. \textbf{H3:} PM is within the threshold of B0. \textbf{H4a:} BPB degrades
monotonically with K/V sharing, and MQA is resolved-worse than B0. \textbf{H4b:} R25 is not resolved from R75 or full RoPE.
\textbf{H5a:} removing value embeddings costs B0 a resolved loss. \textbf{H5b/c:} GQA's gap is smaller without value
embeddings (a positive interaction, probably unresolved). \textbf{H6:} from the Chinchilla data term ($B=410.7$,
$\beta=0.28$, $B/D^{\beta}\approx1.56$ nats at $D=0.44$B), 4.2\% / 6.4\% more tokens should be worth about
$0.018/0.027$ nats, i.e.\ roughly $0.005$--$0.006$ / $0.008$--$0.009$ BPB at 4.2--4.8 bytes per token.
\textbf{HIa--c (inference):} decode never materialises expanded K/V; at batch 1 the variants decode within 5\% of each
other; at batch 64, GQA's decode advantage grows with prompt length up to the training context and then plateaus.

\paragraph{Statistical procedure.}
A difference counts as \emph{resolved} only if it exceeds $2\sqrt{2}\,s$, a two-sigma bound on the difference of two
single runs, where $s$ is the s.d.\ of the three B0 seeds. This rule is crude: $s$ has two degrees of freedom. By
simulation it flags about 18\% of null single-vs-single contrasts, so across our $\approx17$ contrasts two or three
spurious ``resolved'' flags are expected~\cite{bouthillier2021variance}. We therefore pre-registered three secondary
analyses and report where they disagree with the primary rule: (i) a bound that accounts for seed means,
$2s\sqrt{1/n_1+1/n_2}$; (ii) a Student-$t$ 95\% interval using $s_p$, the s.d.\ pooled over the B0 and GQA-2 groups
(df$\,=4$); and (iii) a \emph{paired} analysis that uses the shared data order of same-seed runs (E2 seed pairs and
E6 twins). Validation curves are compared only between runs with the same token target, because a longer run's
learning-rate decay starts later.

\subsection{Inference benchmark}
After training, each final checkpoint is benchmarked on the H100 with fixed random prompts, greedy decoding of exactly 256
tokens, a batch size $\in\{1,4,16,64\}$ and a prompt length $\in\{128, 1792, 8192\}$. A 1{,}792-token prompt plus 256
generated tokens exactly fills the 2{,}048 training context. The 8{,}192-token prompt exceeds the context, so we use it
only to measure speed and memory, with no quality claim. We report prefill throughput, time to first token, decode
throughput (the remaining 255 tokens), peak memory and the allocated KV cache, and compare them with the analytic cache
size and the K/V bytes read per decoded token. Execution is eager. Weights are cast to BF16 once, and the cache position
is kept on the host to avoid a GPU synchronisation per step.

\subsection{Validity checks before the main runs}\label{sec:validity}
CPU unit tests check that full RoPE reproduces the historical code exactly, that the un-rotated dimensions pass through
unchanged, that GQA projection and cache shapes are correct, that cached decoding matches full-sequence decoding, and that
resuming from a checkpoint is bit-identical. On the H100, a preflight for all eight architectures compares FA3 with
PyTorch SDPA (logits, loss and gradients) and the compiled model with eager execution (per-tensor relative gradient
error $\le5\%$) across the sliding-window boundary. The largest observed error was \TBD{x.x\%} on
\TBD{parameter}. A 10-update pilot per architecture then confirmed FA3 selection, no out-of-memory error anywhere in the
inference grid, a decode working set below half the cache (so K/V heads are not expanded), and data-loader wait time
below 2\% of each update, so measured throughput reflects the GPU rather than the CPU packer. All pilots passed and no
setting was changed afterwards (\href{https://github.com/mcxraider/modded-nanochat/blob/main/results/pilot_d8-v5.md}{\texttt{results/pilot\_d8-v5.md}}).
The suite cost \TBD{\$xx.xx} against a projected \$21.38.

% ---------------------------------------------------------------------------
\section{Results}
% ---------------------------------------------------------------------------
\subsection{Noise floor (E1)}
The three B0 seeds reach test BPB \TBD{a}, \TBD{b} and \TBD{c} (mean \TBD{m}), so $s=\TBD{s}$ and the resolution
threshold is $2\sqrt2\,s=\TBD{t}$ BPB. The three GQA-2 seeds give $s_\text{GQA-2}=\TBD{s2}$ and the pooled
$s_p=\TBD{sp}$ (df$\,=4$). The prediction $s\in[0.001,0.003]$ \TBD{held / failed}. \TBD{Comment on whether the two
groups' noise is similar, which the threshold assumes. Comment on validation--test agreement: is the gap between
validation and test BPB constant across runs, i.e.\ is the test shard simply harder/easier?}

\begin{figure}[t]
\centering
\resultfigure{learning_curves_d8.pdf}{0.98\linewidth}{Validation BPB vs.\ update for runs grouped by token target
(1{,}680 and 1{,}750 updates); training loss inset.}
\caption{Learning curves. Curves are compared only within the same token target.}
\label{fig:curves}
\end{figure}

\subsection{Quality at fixed compute (E2--E5)}
\begin{table}[t]
\centering\small
\caption{Final results. $\Delta$ is test BPB minus the primary reference (negative is better). ``Res.'' is the primary
rule $|\Delta|>2\sqrt2 s$, ``means'' the seed-mean rule, and ``$t$'' the pooled-$t$ 95\% interval. Throughput is the
median steady-state training tokens/s. Seed groups are means of three runs.}
\label{tab:results}
\begin{tabular}{@{}lrrrrccc r@{}}
\toprule
Run & Test BPB & Val BPB & Reference & $\Delta$ test & Res. & means & $t$ & Tok/s (rel.\ B0) \\
\midrule
B0 (3 seeds) & \TBD{} & \TBD{} & --- & --- & --- & --- & --- & 1.000 \\
GQA-2 (3 seeds) & \TBD{} & \TBD{} & B0 & \TBD{} & \TBD{} & \TBD{} & \TBD{} & \TBD{} \\
GQA-2-R75 & \TBD{} & \TBD{} & GQA-2 & \TBD{} & \TBD{} & \TBD{} & \TBD{} & \TBD{} \\
GQA-2-R25 & \TBD{} & \TBD{} & GQA-2 & \TBD{} & \TBD{} & \TBD{} & \TBD{} & \TBD{} \\
GQA-2-PM & \TBD{} & \TBD{} & B0 & \TBD{} & \TBD{} & \TBD{} & \TBD{} & \TBD{} \\
MQA & \TBD{} & \TBD{} & B0 & \TBD{} & \TBD{} & \TBD{} & \TBD{} & \TBD{} \\
B0-noVE & \TBD{} & \TBD{} & B0 & \TBD{} & \TBD{} & \TBD{} & \TBD{} & \TBD{} \\
GQA-2-noVE & \TBD{} & \TBD{} & B0-noVE & \TBD{} & \TBD{} & \TBD{} & \TBD{} & \TBD{} \\
GQA-2-tok & \TBD{} & \TBD{} & B0 & \TBD{} & \TBD{} & \TBD{} & \TBD{} & \TBD{} \\
MQA-tok & \TBD{} & \TBD{} & B0 & \TBD{} & \TBD{} & \TBD{} & \TBD{} & \TBD{} \\
\bottomrule
\end{tabular}
\end{table}

\paragraph{GQA-2 vs.\ MHA (H2a).}
At fixed compute, GQA-2 reaches \TBD{x} test BPB against B0's \TBD{y} ($\Delta=\TBD{d}$, \TBD{resolved / unresolved}).
The paired analysis over the three same-seed pairs gives a mean difference of \TBD{d} with 95\% interval \TBD{[l,u]}, and
$s_d/(\sqrt2 s)=\TBD{r}$, so the shared data order removes \TBD{\%} of the seed noise. \TBD{Interpretation: does H2a hold?
Note that this contrast changes K/V heads, value-embedding size and token count at once; E5 and E6 below decompose it.}

\paragraph{K/V dose--response (H4a).}
\TBD{BPB at $n_\text{kv}=4,2,1$; monotone or not? MQA vs B0 resolved? Recall that MQA has 6.4\% more tokens and a
quarter of B0's value-embedding parameters, and that its value gate is a $1\times32$ matrix, for which Muon's
orthogonalised update reduces to a normalised vector.}

\paragraph{Partial RoPE (H2c, H4b).}
\TBD{R75 and R25 against GQA-2, all at 1{,}750 updates, so curves are directly comparable (Fig.~\ref{fig:dose}).
Interpret with the frequency argument from \S2: R25 keeps the extreme frequencies but drops 75\% of the rotated
channels in favour of position-free ones.}

\begin{figure}[t]
\centering
\resultfigure{dose_response.pdf}{0.49\linewidth}{Test BPB vs.\ K/V heads (4, 2, 1) and vs.\ rotary fraction (1, 0.75, 0.25), with the $\pm2\sqrt2 s$ band.}
\hfill
\resultfigure{budget_comparison.pdf}{0.49\linewidth}{E6: test BPB vs.\ estimated FLOPs; arrows from each matched-token ablation to its fixed-compute twin; budget $C$ marked.}
\caption{Left: dose--response along the K/V and rotary axes. Right: what the extra tokens buy (E6).}
\label{fig:dose}
\end{figure}

\paragraph{Value embeddings (E5, H5a--c).}
\begin{table}[t]
\centering\small
\caption{E5 $2\times2$ on test BPB. The interaction is $(\text{GQA-2}-\text{B0})-(\text{GQA-2-noVE}-\text{B0-noVE})$,
with two-sigma bound $2s\sqrt{1/3+1/3+1+1}\approx3.3s$.}
\label{tab:ve}
\begin{tabular}{@{}lccc@{}}
\toprule
& Value emb.\ on & Value emb.\ off & Effect of removing \\
\midrule
MHA ($n_\text{kv}=4$) & \TBD{} (B0, $n=3$) & \TBD{} & \TBD{} \\
GQA-2 ($n_\text{kv}=2$) & \TBD{} ($n=3$) & \TBD{} & \TBD{} \\
Effect of GQA-2 & \TBD{} & \TBD{} & interaction \TBD{} (bound \TBD{}) \\
\bottomrule
\end{tabular}
\end{table}
\TBD{Is the value-embedding effect resolved (H5a)? Is the GQA gap smaller without value embeddings (H5b)? How large is the
interaction relative to its bound? If the GQA penalty largely disappears without value embeddings, then the ``GQA cost''
in this architecture is mostly a value-embedding capacity cost, which matters for anyone porting GQA into a model with
value residuals.}

\paragraph{Spending the saving on parameters instead of tokens (E3, H3).}
\TBD{GQA-2-PM vs B0 and vs GQA-2. PM and GQA-2 are two ways to spend the same FLOP saving: on 256 extra MLP units or on 70
extra updates. Which wins? Remember that PM still has 33.5M fewer embedding parameters than B0, and that
Muon gives \texttt{c\_fc} about 6\% more learning rate in the wider MLP.}

\subsection{What the extra tokens buy (E6, H6)}
\TBD{GQA-2(s1) $-$ GQA-2-tok $=\Delta_1$ and MQA $-$ MQA-tok $=\Delta_2$ against the predicted $-0.005$ and $-0.009$ BPB.
GQA-2-tok $-$ B0 and MQA-tok $-$ B0 are single-seed estimates of the per-token cost of K/V sharing. Decompose the E2 result:
fixed-compute gap $=$ per-token cost $+$ data gain. State the paired-noise assumption ($2s_d$ threshold) and whether the
Chinchilla-based prediction was of the right order.}

\subsection{Training throughput and memory}\label{sec:speed}
The pilots (8 timed updates each) and full runs agree that GQA buys little training speed at this size
(Table~\ref{tab:speed}). GQA-2 removes 4.2\% of estimated FLOPs but is only \TBD{1.6}\% faster, and MQA removes 6.4\%
and is \TBD{5.5}\% faster. Throughput in FLOP/s falls slightly as K/V heads are removed (379 to 369 TFLOP/s for GQA-2 in
the pilot), which is consistent with the narrower $512\times256$ K/V matmuls using the tensor cores less efficiently than
the square ones and with attention, which GQA does not shrink, taking a larger share of the step. Two observations show
the limits of the FLOP budget. First, removing the value embeddings leaves the FLOP estimate unchanged but makes B0
\TBD{2.0}\% faster: the lookups, gates and their AdamW state cost real time that the estimator does not count. Second,
GQA-2-PM is FLOP-matched to B0 but \TBD{1.5}\% slower, because its 2{,}304-wide MLP is less efficient on the hardware.
Under a wall-clock budget, both effects would shift tokens towards the variants without value embeddings and away from PM.
Partial RoPE changes throughput by under 1\% (H2c). \TBD{Peak training VRAM per architecture; evaluation peak separately.}

\begin{table}[t]
\centering\small
\caption{Training throughput (full-run steady-state median; pilot value in parentheses) and achieved TFLOP/s ($=$
tokens/s $\times$ FLOPs/token; H100 dense BF16 peak 989). \TBD{Replace pilot-only entries with full-run values.}}
\label{tab:speed}
\begin{tabular}{@{}lrrrr@{}}
\toprule
Architecture & Tokens/s (M) & vs.\ B0 & FLOP saving & Achieved TFLOP/s (MFU) \\
\midrule
B0 & \TBD{} (1.205) & 1.000 & --- & \TBD{} (379, 38\%) \\
GQA-2 & \TBD{} (1.223) & \TBD{} (1.016) & 4.2\% & \TBD{} (369, 37\%) \\
GQA-2-R75 & \TBD{} (1.231) & \TBD{} (1.022) & 4.2\% & \TBD{} (372, 38\%) \\
GQA-2-R25 & \TBD{} (1.230) & \TBD{} (1.021) & 4.2\% & \TBD{} (371, 38\%) \\
GQA-2-PM & \TBD{} (1.187) & \TBD{} (0.985) & 0.0\% & \TBD{} (373, 38\%) \\
MQA & \TBD{} (1.271) & \TBD{} (1.055) & 6.4\% & \TBD{} (376, 38\%) \\
B0-noVE & \TBD{} (1.228) & \TBD{} (1.020) & 0.0\% & \TBD{} (386, 39\%) \\
GQA-2-noVE & \TBD{} (1.242) & \TBD{} (1.031) & 4.2\% & \TBD{} (375, 38\%) \\
\bottomrule
\end{tabular}
\end{table}

\subsection{Inference: where the K/V saving appears}
We expected the K/V saving to matter only when K/V reads dominate the bytes moved per decode step. Analytically, at
batch 1 and a 128-token prompt, the 84\,MB of BF16 matmul weights dominate, and K/V is only 5\% of the bytes read per step in
B0, so the ideal bandwidth-bound speedup of GQA-2 is $1.08\times$. At batch 64 and the full 2{,}048 positions, K/V is
94\% of the traffic, so the ideal speedups are $1.9\times$ (GQA-2) and $3.4\times$ (MQA). With the window
capping reads at 1{,}025/2{,}049 positions, they barely change from the 1{,}792- to the 8{,}192-token prompt, while the
allocated cache grows $4\times$ (8.25\,GB for B0 at batch 64).
\TBD{Measured results (Fig.~\ref{fig:inference}): HIa (decode working set vs.\ half the cache), HIb (batch-1 spread), HIc (batch-64
speedups vs.\ the ideal ratios above). Explain the gap between measured and ideal, e.g.\ eager kernel-launch overhead at
small batch, or FA3's decode kernel not being purely bandwidth-bound.}

\begin{figure}[t]
\centering
\resultfigure{inference_comparison_d8.pdf}{0.98\linewidth}{Decode tokens/s, TTFT, peak VRAM and allocated KV cache vs.\ prompt length, per batch size, for B0, GQA-2 and MQA.}
\caption{Inference benchmark on one H100 (eager, BF16, greedy, 256 new tokens).}
\label{fig:inference}
\end{figure}

% ---------------------------------------------------------------------------
\section{Discussion}
% ---------------------------------------------------------------------------
\paragraph{What the experiments establish.}
\TBD{Two or three sentences: the decomposition of GQA's fixed-compute result into per-token cost, data gain and
value-embedding capacity; partial RoPE; where inference gains appear. State which hypotheses held and which failed.}

\paragraph{Predictions vs.\ outcomes.}
\begin{table}[h]
\centering\small
\begin{tabularx}{\linewidth}{@{}lXl@{}}
\toprule
Hyp. & Prediction (registered before results) & Outcome \\
\midrule
E1 & $s\in[0.001,0.003]$ BPB & \TBD{} \\
H2a & GQA-2 not resolved-worse than B0 at fixed compute & \TBD{} \\
H2b & GQA-2 training-throughput gain $<5\%$ & \TBD{held (pilot: 1.6\%)} \\
H2c & R75 unresolved in BPB; $<2\%$ throughput difference & \TBD{} \\
H3 & PM within threshold of B0 & \TBD{} \\
H4a & Monotone degradation with K/V sharing; MQA resolved-worse & \TBD{} \\
H4b & R25 unresolved from R75/full & \TBD{} \\
H5a & Removing value embeddings costs B0 a resolved loss & \TBD{} \\
H5b/c & GQA gap smaller without value embeddings; interaction unresolved & \TBD{} \\
H6 & Extra tokens worth $\approx0.005$ / $0.009$ BPB (GQA-2 / MQA) & \TBD{} \\
HIa & No K/V-sized temporary in decode & \TBD{held in all pilots} \\
HIb & Batch-1 decode within 5\% across variants & \TBD{} \\
HIc & Batch-64 GQA speedup grows to 1{,}792 tokens, then plateaus & \TBD{} \\
\bottomrule
\end{tabularx}
\end{table}

\paragraph{Limitations and threats to validity.}\label{sec:limits}
\begin{itemize}
\item \emph{Statistical power.} Three seeds give $s$ two degrees of freedom. Most variants are single runs, and no
  multiplicity correction is applied across $\approx17$ contrasts. A true effect of $2s$ against a three-seed mean
  is detected only about 40\% of the time. ``Unresolved'' therefore means ``not distinguishable at this budget'', not
  ``equal''.
\item \emph{The FLOP measure.} The estimator ignores embedding lookups, norms, softmax and rotary work, so value embeddings
  are free and partial RoPE saves nothing by its measure. It also charges every token the full attention window rather than the
  causal average, which slightly understates GQA's relative saving. Measured throughput (\S\ref{sec:speed}) shows how
  this differs from a time budget.
\item \emph{Architectural confounds.} The value embeddings shrink with $n_\text{kv}$ (E5 addresses this only at
  $n_\text{kv}\in\{4,2\}$). The value gate becomes a single row at MQA. NorMuon's factored second moment is reduced along
  the longer axis, so the square MHA K/V matrices receive per-row statistics while the GQA ones receive per-column
  statistics. PM's wider \texttt{c\_fc} gets about 6\% more Muon learning rate. None of these is isolated.
\item \emph{No per-variant tuning.} All hyperparameters were tuned upstream for MHA. Part of any GQA penalty could come
  from this mismatch, which a learning-rate sweep at $0.5\times$ and $2\times$ would test.
\item \emph{Scale and duration.} All results are for a 42M-scaling-parameter model trained on about 0.45B tokens. GQA's quality
  cost is known to shrink with model size~\cite{ainslie2023gqa}, so we make no claim about transfer to larger models.
\item \emph{Metric scope.} BPB on ClimbMix measures in-distribution language modelling only, not downstream tasks or the
  use of long context, which is where positional encoding choices are expected to matter most.
\end{itemize}

\paragraph{Future work.}
A d12 replication to test whether the per-token cost of GQA shrinks with size; a learning-rate sweep for B0 and GQA-2;
fixing the value-embedding width to $n_\text{head}\cdot d_\text{head}$ independently of $n_\text{kv}$, which would
remove the main confound; and a wall-clock-budget variant of the suite.

% ---------------------------------------------------------------------------
\section{Conclusion}
% ---------------------------------------------------------------------------
\TBD{Three to four sentences answering the research question with the measured numbers.}

\paragraph{Reproducibility.}
Every run is defined in one table (\texttt{experiments/config.py}), from which all configs are generated. The whole suite
runs with a single Modal command, and the tables and figures in this report are regenerated without a GPU from the
saved results in \texttt{results/} (see \texttt{readme.md}). Each run records its config, code commit and diff,
hardware, software versions, data and tokenizer hashes, and data order.

{\footnotesize
\setlength{\bibsep}{1pt}
\bibliographystyle{plainnat}
\bibliography{references}
}

\end{document}
